\subsection{阿廷环上的投射模}

命题 4. —— 左阿廷环上的每个模都具有一个投射覆盖。

设 A 是一个左阿廷环，M 是 A 上的一个模。A 的根 r 是一个幂零双边理想，且环 A/r 是半单的（VIII, 第 173 页，命题 1）。由此推出 A/r-模 M/rM 是投射的。于是该断言由 VIII, 第 164 页的命题 11 得出。

命题 5. —— 设 A 是一个左阿廷环，r 是它的根。

\begin{enumerate}
\def\labelenumi{\alph{enumi})}
\item
  设 P 是一个投射 A-模。用 u 表示从 P 到 P/rP 的典范映射。则 \((P,u)\) 是 P/rP 的一个投射覆盖。特别地，A-模 P 是有限生成的当且仅当 P/rP 是有限生成的。
\item
  设 M 是环 A/r 上的一个模。则把 M 视为 A-模时，它具有一个投射覆盖。若 \((P,u)\) 是这样一个投射覆盖，则取商时 u 诱导出从 P/rP 到 M 的一个同构。此外，P 是不可分解的当且仅当 M 是单的。
\item
  设 M 与 M' 是环 A/r 上的模。设 \((P, u)\) 与 \((P', u')\) 分别是 A-模 M 与 M' 的投射覆盖。则 M 与 M' 同构当且仅当 P 与 P' 同构。
\end{enumerate}

A 的根 r 是一个幂零双边理想（VIII, 第 173 页，命题 1）。于是断言 a) 由 VIII, 第 163 页的推论 1 及 VIII, 第 161 页的注 2 得出。

我们来证明断言 b)。M 的投射覆盖 \((P, u)\) 的存在性由命题 4 及 VIII, 第 162 页命题 9 的从 P/rP 到 M 的同构得出。由于环 A/r 是半单的，这个环上的模是不可分解的当且仅当它是单的。最后的断言由 VIII, 第 160 页的推论 2 得出。

我们来证明 c)。设 f 是从 P 到 \(P'\) 的一个同构；它诱导出一个同构 \(\overline{f}\)：从 \(P/rP\) 到 \(P'/rP'\)。由 b)，这些商分别同构于 M 与 \(M'\)。因此 M 同构于 \(M'\)。反之，由 VIII, 第 161 页的命题 8，从 M 到 \(M'\) 的每个同构 f 都可提升为一个同构 \(\widetilde{f}\)（从 P 到 \(P'\)），使得 \(f \circ u = u' \circ \widetilde{f}\)。

推论. —— 对投射 A-模的每个类 P，对应环 A/r 上模 P/rP 的类 \(\varphi(\mathrm{P})\)。环 A/r 上的每个模类（相应地，有限生成模类）都具有形式 \(\varphi(\mathrm{P})\)，其中 P 是唯一确定的投射 A-模类（相应地，有限生成投射 A-模类）。

设 A 是一个左阿廷环，\(\mathfrak{r}\) 是它的根。用 \(\mathcal{S}\) 表示单 A-模的类所成之集；它是有限的（VIII, 第 51 页）。由 VIII, 第 174 页的命题 2，它可以典范地等同于半单环 \(\mathrm{A} / \mathfrak{r}\) 上单模的类所成之集。对每个 \(\lambda \in \mathcal{S}\)，设 \(\mathrm{S}_{\lambda}\) 是类为 \(\lambda\) 的一个单 A-模，并选取一个投射覆盖 \((\mathrm{P}_{\lambda}, u_{\lambda})\)（它是 \(\mathrm{S}_{\lambda}\) 的）（VIII, 第 175 页，命题 4）。由命题 5，A-模 \(\mathrm{P}_{\lambda}\) 是投射的且不可分解的，且 \(u_{\lambda}\) 定义了从 \(\mathrm{P}_{\lambda} / \mathfrak{r}\mathrm{P}_{\lambda}\) 到 \(\mathrm{S}_{\lambda}\) 的一个同构。此外，若 P 是一个投射且不可分解的 A-模，则存在唯一的 \(\lambda \in \mathcal{S}\) 使得 P 同构于 \(\mathrm{P}_{\lambda}\)：它就是唯一的 \(\lambda \in \mathcal{S}\)，使 \(\mathrm{P} / \mathfrak{r}\mathrm{P}\) 同构于 \(\mathrm{S}_{\lambda}\)。

命题 6. —— 设 A 是一个左阿廷环，r 是它的根。设 P 是一个投射 A-模。

\begin{enumerate}
\def\labelenumi{\alph{enumi})}
\item
  A-模 \(\overline{\mathrm{P}} = \mathrm{P} / \mathfrak{r}\mathrm{P}\) 是半单的，且 A-模 \(\mathrm{P}\) 同构于 \(\bigoplus_{\lambda \in \mathcal{S}}\mathrm{P}_{\lambda}^{([\overline{\mathrm{P}}:\lambda ])}\)。
\item
  设 \((\mathrm{Q}_{i})_{i\in\mathrm{I}}\) 是 P 的一个以 P 为直和的不可分解投射子模族。则对每个 \(\lambda\in S\)，集 \(\mathrm{I}(\lambda)\)（由 \(i\in I\) 中使 \(Q_{i}\) 同构于 \(P_{\lambda}\) 者组成）的基数等于 \([\overline{P}:\lambda]\)。
\end{enumerate}

\(\overline{P}\) 是半单的这一事实由命题 2（VIII, 第 174 页）得出。A-模 \(Q = \oplus_{\lambda \in S} P_{\lambda}^{([\overline{P}: \lambda])}\) 是投射的。由于 \(P_{\lambda}/rP_{\lambda}\) 同构于 \(S_{\lambda}\)，商 Q/rQ 同构于 \(\oplus_{\lambda \in S} S_{\lambda}^{([\overline{P}: \lambda])}\)，即同构于 \(\overline{P} = P/rP\)。由命题 5，A-模 P 与 Q 同构。

设给定了 P 的一个以 P 为直和的投射且不可分解的子模族 \((\mathrm{Q}_{i})_{i\in\mathrm{I}}\)。由命题 5，A-模 \(Q_{i}/rQ_{i}\) 对每个 \(i\in I\) 都是单的。对 \(\lambda\in S\)，用 \(\mathrm{I}(\lambda)\) 表示由 \(i\in I\)（使 \(Q_{i}/rQ_{i}\) 同构于 \(S_{\lambda}\)、从而同构于 \(P_{\lambda}/rP_{\lambda}\)）组成的集；它也是由 \(i\in I\)（使 \(Q_{i}\) 同构于 \(P_{\lambda}\)）组成的集。由于 \(\overline{P}=P/rP\) 是族 \((\mathrm{Q}_{i}/rQ_{i})_{i\in\mathrm{I}}\) 的直和，由 VIII, 第 32 页的定理 1，我们有 \(\operatorname{Card}(\mathrm{I}(\lambda))=[\overline{\mathrm{P}}:\lambda]\)。

例. —— 取 P 等于 \(A_{s}\)。对 \(\lambda \in S\)，将单 A-模 \(S_{\lambda}\) 的交换子之反体记作 \(D_{\lambda}\)，并将 \(S_{\lambda}\) 视为体 \(D_{\lambda}\) 上右向量空间时的维数记作 \(m(\lambda)\)。我们知道 \(m(\lambda)\) 等于重数 \([A_{s}/rA_{s} : \lambda]\)（VIII, 第 143 页，命题 11）。于是，A-模 \(A_{s}\) 同构于 \(\oplus_{\lambda \in S} P_{\lambda}^{m(\lambda)}\)。

